% section: 解を直接求めず,近づく数列を設計する

\section{解を直接求めず,近づく数列を設計する}

\subsection{テイラー展開は何をしているのか}

関数 $f(x)$ を $x=a$ の近くで調べる.
接線による近似は,
\[
  f(x)\approx f(a)+f'(a)(x-a)
\]
である.
二次の項まで含めれば,
\[
  f(x)
  \approx f(a)+f'(a)(x-a)
  +\frac{f''(a)}{2}(x-a)^2
\]
となる.

テイラー展開は,関数を多項式で近似する道具である.
ただし,「近似だから適当でよい」という意味ではない.
どの点の近くで,何次まで使うかを決めることで,関数の局所的な情報を数値計算に変換している.

\subsection{ニュートン法の導出}

方程式
\[
  f(x)=0
\]
の解を探す.
現在の近似値を $x_n$ とし,$x_n$ の近くでは関数を接線で置き換える.
\[
  0\approx f(x_n)+f'(x_n)(x-x_n)
\]
を $x$ について解くと,
\[
  x=x_n-\frac{f(x_n)}{f'(x_n)}
\]
である.
この値を次の近似値に採用する.
\[
  x_{n+1}=x_n-\frac{f(x_n)}{f'(x_n)}
\]

ニュートン法は,方程式を解く公式ではない.
方程式を,収束することを期待した漸化式へ変換する設計法である.

\subsection{平方根の例と二次収束}

\[
  f(x)=x^2-2
\]
にニュートン法を使うと,
\[
  x_{n+1}=\frac{x_n^2+2}{2x_n}
\]
となる.
$x_0=1$ とすれば,
\[
  x_1=\frac32,\quad
  x_2=\frac{17}{12},\quad
  x_3=\frac{577}{408}
\]
である.

$\alpha=\sqrt2$ とおくと,
\[
  \begin{aligned}
    x_{n+1}-\alpha
     & =\frac{x_n^2+2}{2x_n}-\alpha             \\
     & =\frac{x_n^2-2\alpha x_n+\alpha^2}{2x_n} \\
     & =\frac{(x_n-\alpha)^2}{2x_n}.
  \end{aligned}
\]

誤差を $e_n=x_n-\alpha$ と書けば,
\[
  e_{n+1}=\frac{e_n^2}{2x_n}
\]
である.
誤差が二乗されるため,十分近くまで来ると正しい桁数が急速に増える.
この性質を二次収束という.

\subsection{一般の単純な根の近くで起こること}

$f(\alpha)=0$ かつ $f'(\alpha)\neq0$ とする.
$x_n=\alpha+e_n$ とおき,二次までテイラー展開すると,
\[
  f(x_n)=f'(\alpha)e_n+\frac{f''(\alpha)}2e_n^2+\cdots
\]
である.
同様に,
\[
  f'(x_n)=f'(\alpha)+f''(\alpha)e_n+\cdots
\]
となる.
これをニュートン法へ代入すると,一次の誤差が打ち消され,

\[
  e_{n+1}=C e_n^2+\text{三次以上の項}
\]
という形になる.
これがニュートン法が,根の近くで二次収束する理由である.

ただし,根が重なっていて $f'(\alpha)=0$ になる場合や,初期値が遠い場合には,この説明はそのまま使えない.
ニュートン法を適用する前に,何を近似しているのか,どの範囲で近似が有効なのかを確認する必要がある.

\subsection{微分を使わないニュートン法}

微分が計算しにくい場合,二つの近似値を使って接線の傾きを近似することができる.
\[
  f'(x_n)\approx
  \frac{f(x_n)-f(x_{n-1})}{x_n-x_{n-1}}
\]
と置くと,
\[
  x_{n+1}
  =
  x_n
  -
  f(x_n)
  \frac{x_n-x_{n-1}}{f(x_n)-f(x_{n-1})}
\]
を得る.
これは割線法と呼ばれる.

ニュートン法と割線法の違いは,単なる公式の違いではない.
ニュートン法は関数の傾きを正確に使い,割線法は過去の値から傾きを推測する.
どちらも「次の値をどう設計すれば,解へ近づくか」という同じ問題への答えである.
